\section{Conclusion}
\label{sec:conclusion}
    In this paper, we presented the design and implementation of a remote train control system, leveraging the capabilities of the deterministic execution model of Lingua Franca. Our architecture integrates real-time video streaming and control command transmission, demonstrating the feasibility of using QUIC-based WebTransport to meet the network induced latency requirements over real-world scenarios.

    Future work will focus on several directions. First, we plan to enhance the operator interface by showing situational-awareness features. Second, we intend to conduct larger-scale field tests to evaluate system performance under more challenging and diverse network conditions, including high-latency, lossy, and intermittent links. Third, we plan to explore multi-train operation scenarios—including conflict detection and resolution—and leverage Federated Lingua Franca for distributed execution across multiple simultaneous train deployments.